% Lineare Funktionen · Funktionswert · Prüfungs-Fokus MSA – bank/lineare-funktionen/e3.jsonl (zusammenbau v0.6, Prüfungs-Fokus)
\einheitenkopf[e3]{Einheit 3 · Punkte und Werte}
\zweigzeile{ab Kl. 8 · P10 ×12}
\setcounter{aufgabe}{0}

% Anlauf (ohne Prüfkennung)
\begin{aufgabe}{Funktionswert – Anlauf}
\begin{teile}
\teil $f(4)$? Berechne den Funktionswert von $f(x) = 2x + 5$ an der Stelle $x = 4$. f(4) = \leerfeld
\teil $f(-3)$? Berechne den Funktionswert von $f(x) = 4x - 1$ an der Stelle $x = -3$. f(-3) = \leerfeld
\end{teile}
\end{aufgabe}

\begin{aufgabe}{Funktionswert – Prüfungsaufgaben}
\begin{teile}
\teil Graph und Nullstelle: Zeichne den Graphen der Funktion f mit $f(x) = -4x + 6$ und berechne die Nullstelle von f. \hfill (P22) \\

\begin{ksys}[xmin=-4,xmax=4,ymin=-4,ymax=7] \end{ksys}
x = \leerfeld
\teil Graph und Nullstelle: Zeichne den Graphen der Funktion f mit $f(x) = -6x + 3$ und berechne die Nullstelle von f. \hfill (P22) \\

\begin{ksys}[xmin=-4,xmax=4,ymin=-4,ymax=4] \end{ksys}
x = \leerfeld
\teil Wann ist der Tank leer? Ein Wassertank enthält 60 Liter. Die Gleichung $y = -0{,}4x + 60$ gibt an, wie viele Liter y nach x Minuten noch im Tank sind. Berechne, nach wie vielen Minuten der Tank leer ist. \hfill (P21) \\ \leerfeld[min]
\teil Wann ist der Akku leer? Ein voll geladener Handyakku verliert beim Spielen 0{,}8 Prozentpunkte je Minute. Die Gleichung $y = -0{,}8x + 100$ gibt den Ladestand y in Prozent nach x Minuten an. Berechne, nach wie vielen Minuten der Akku leer ist. \hfill (P21) \\ \leerfeld[min]
\teil Nullstelle ablesen: Gib die Nullstelle der Funktion g an. \hfill (P19) \\

\begin{ksys}[xmin=-5,xmax=5,ymin=-5,ymax=5,ablesen] \gerade{1.5}{3}{g} \end{ksys}
x = \leerfeld
\teil Nullstelle ablesen: Gib die Nullstelle der Funktion g an. \hfill (P19) \\

\begin{ksys}[xmin=-5,xmax=5,ymin=-5,ymax=5,ablesen] \gerade{-0.5}{2}{g} \end{ksys}
x = \leerfeld
\teil Gerade und Parabel: Zeichne die Gerade f mit $f(x) = 3x - 3$ in das Koordinatensystem, gib die Koordinaten des sichtbaren Schnittpunkts von f mit der Parabel p an und entscheide, ob sich f und p in zwei Punkten schneiden. \hfill (P23) \\

\begin{ksys}[xmin=-5,xmax=5,ymin=-6,ymax=8] \parabel{-1}{1}{4}{p} \end{ksys}
S( \leerfeld | \leerfeld ) \quad \janein
\teil Gerade und Parabel: Zeichne die Gerade f mit $f(x) = 4x + 2$ in das Koordinatensystem, gib die Koordinaten des sichtbaren Schnittpunkts von f mit der Parabel p an und entscheide, ob sich f und p in zwei Punkten schneiden. \hfill (P23) \\

\begin{ksys}[xmin=-5,xmax=5,ymin=-6,ymax=8] \parabel{-1}{0}{7}{p} \end{ksys}
S( \leerfeld | \leerfeld ) \quad \janein
\teil y-Koordinate: Der Punkt A liegt auf der Geraden $y = \frac{1}{3}x - 2$ und hat die x-Koordinate $-9$. Berechne die y-Koordinate von A. \hfill (P17) \\ y = \leerfeld
\teil y-Koordinate: Der Punkt A liegt auf der Geraden $y = \frac{3}{4}x + 2$ und hat die x-Koordinate $-8$. Berechne die y-Koordinate von A. \hfill (P17) \\ y = \leerfeld
\teil Liegt P auf der Geraden? Prüfe rechnerisch, ob $P(-4|5)$ auf der Geraden $f(x) = -\frac{1}{2}x + 3$ liegt. \hfill (P26F) \\ \janein
\teil Liegt P auf der Geraden? Prüfe rechnerisch, ob $P(-2|4)$ auf der Geraden $f(x) = -\frac{3}{2}x - 1$ liegt. \hfill (P26F) \\ \janein
\teil Welche zwei Aussagen stimmen? Gegeben ist $f(x) = 2x + 6$. Kreuze die beiden richtigen Aussagen an. \hfill (P21) \\ \kreuz{$Q(6|0)$ liegt auf der Geraden f.} \\ \kreuz{Die Gerade f fällt.} \\ \kreuz{Die Gerade f schneidet die y-Achse in $P(0|6)$.} \\ \kreuz{Die Gerade f ist parallel zur Geraden $y = 2x$.}
\teil Welche zwei Aussagen stimmen? Gegeben ist $f(x) = -3x + 9$. Kreuze die beiden richtigen Aussagen an. \hfill (P21) \\ \kreuz{$Q(-3|0)$ liegt auf der Geraden f.} \\ \kreuz{Die Gerade f fällt.} \\ \kreuz{Die Gerade f schneidet die y-Achse in $P(0|-3)$.} \\ \kreuz{Der Punkt $R(2|3)$ liegt auf der Geraden f.}
\teil Welcher Punkt liegt nicht auf der Geraden? Gegeben ist $f(x) = -6x + 2$. Kreuze den Punkt an, der nicht auf der Geraden f liegt. \hfill (P23) \\ \kreuz{$(-1|8)$} \\ \kreuz{$(2|-10)$} \\ \kreuz{$(-3|16)$}
\teil Welcher Punkt liegt nicht auf der Geraden? Gegeben ist $f(x) = -5x + 4$. Kreuze den Punkt an, der nicht auf der Geraden f liegt. \hfill (P23) \\ \kreuz{$(1|-1)$} \\ \kreuz{$(-2|-6)$} \\ \kreuz{$(3|-11)$}
\end{teile}
\end{aufgabe}

\medskip
\uebersichtskasten{Funktionswert: x einsetzen. \quad f(x) = 2x + 1, x = 3: f(3) = 2 · 3 + 1 = 7 \\ Punktprobe: x einsetzen, mit y vergleichen. \quad P(3 | 7) liegt auf f, weil f(3) = 7 (wA) \\ Nullstelle: f(x) = 0 setzen, nach x auflösen.  0 = 2x + 1 → x = −0,5}
